\documentclass[11pt,a4paper]{letter}
\usepackage[utf8]{inputenc}
\usepackage[T1]{fontenc}
\usepackage[margin=2.5cm]{geometry}
\usepackage{hyperref}

\signature{Luiz Diego Vidal Santos\\
Programa de P\'{o}s-Gradua\c{c}\~{a}o em Planejamento Territorial\\
Universidade Estadual de Feira de Santana\\
Feira de Santana, BA 44036-900, Brazil\\
\texttt{ldvsantos@uefs.br}\\
ORCID: 0000-0001-8659-8557}

\date{\today}

\begin{document}

\begin{letter}{The Editor-in-Chief\\
\textit{Earth Critical Zone}\\
Elsevier}

\opening{Dear Editor,}

We are pleased to submit the enclosed manuscript entitled \textbf{``Parametric identifiability of a multiplicative USLE erodibility correction for hydraulically impeded Plinthosols''} for consideration as an Original Research Article in \textit{Earth Critical Zone}.

The erodibility factor ($K$) of the Universal Soil Loss Equation remains the most widely used metric of soil susceptibility to water erosion, yet its estimation via the standard nomograph is restricted to surface-layer attributes and systematically underestimates erosive susceptibility in tropical soils governed by subsurface processes. Plinthosols, which are widespread across tropical landscapes, combine subsurface hydraulic impedance, aluminum phytotoxicity and geotechnical vulnerability --- three convergent mechanisms absent from the original $K$-factor formulation that operate within the critical zone where water, soil and biota interact.

This study proposes and computationally evaluates a multiplicative correction factor ($K_{\mathrm{plint}}$) that translates these three mechanisms as independent amplifiers of $K_{\mathrm{RUSLE}}$, with five free parameters. Through five computational tests anchored in field-measured properties of a Haplic Plinthosol from the Coastal Tablelands of Sergipe (northeastern Brazil), complemented by virtual profiles of a Ferralsol and an Acrisol, we demonstrate that the three hydraulic-biological parameters are recoverable with bias below 15\% even under 30\% observational noise, while Sobol global sensitivity analysis confirms that subsurface hydraulic impedance governs 63--78\% of amplification variance in the Plinthosol. The Ferralsol negative control consistently yields amplification near unity, validating the multi-class experimental design, and Green--Ampt infiltration modeling coupled with simplified WEPP-type erosion simulation preserves the pedological hierarchy across soil classes.

We believe this manuscript aligns closely with the scope of \textit{Earth Critical Zone}, particularly its focus on the interactions among water, soil, air, organisms and human activities within the critical zone, and the development of monitoring and assessment tools for soil degradation in data-scarce tropical environments. The integration of hydraulic, geochemical and geotechnical processes into a single erodibility correction factor addresses the interdisciplinary nature that \textit{Earth Critical Zone} promotes.

The simulation scripts, input parameters and generated figures are publicly available at Zenodo (\url{https://doi.org/10.5281/zenodo.19483403}).

This manuscript has not been published previously, is not under consideration for publication elsewhere, and all authors have approved the submitted version. There are no competing interests to declare.

\closing{Sincerely,}

\end{letter}
\end{document}
